
%\subsection{Steenrod squares}
%
%We work in the derived category \(\cD(\Ftwo)\) of unbounded chain complexes of vector spaces over the field with two elements.
%For the structural identifications below, we assume that \(V\) is \defn{homologically finite}, that is, \(H(V)\) is finite-dimensional over \(\Ftwo\).
%Equivalently, \(V\) is a finite direct sum of shifts of \(\Ftwo\) in \(\cD(\Ftwo)\).
%More generally, the same statements hold under any hypotheses guaranteeing convergence of the relevant Tate spectral sequences and excluding completion issues.
%
%\subsubsection{Shifted Tate sequence}
%
%We use the symmetry of the tensor product to define the functors
%\[
%V \mapsto (V^{\ot 2})_{h\cyc_2}
%\qquad\text{and}\qquad
%V \mapsto (V^{\ot 2})^{h\cyc_2},
%\]
%where \((-)_{h\cyc_2}\) and \((-)^{h\cyc_2}\) are respectively the functors of homotopy orbits and homotopy fixed points.
%The homotopy fibre of the norm map defines another functor
%\[
%T_2 V \defeq \fib\Big((V^{\ot 2})_{h\cyc_2} \xrightarrow{\mathrm{Nm}} (V^{\ot 2})^{h\cyc_2}\Big),
%\]
%naturally equivalent to the shifted Tate construction \(T_2 V \simeq (V^{\ot 2})^{t\cyc_2}[-1]\), and which fits in the \textit{shifted Tate sequence}
%\begin{equation}\label{eq:tate_sequence}
%	(V^{\ot 2})^{h\cyc_2}[-1] \to T_2 V \to (V^{\ot 2})_{h\cyc_2}.
%\end{equation}
%
%With this convention, for every \(j \in \Z\),
%\[
%H_j(T_2 \Ftwo) \cong \Ftwo\set{\hat x_j},
%\]
%so \(H(T_2 \Ftwo)\) is a free one-dimensional module over the ring \(\Ftwo[u,u^{-1}]\), where \(u\) acts by
%\[
%u \, \hat x_j = \hat x_{j + 1}.
%\]
%
%\subsubsection{Tate spectral sequence}
%
%Let \(M\) be a chain complex with an action of \(\cyc_2\), and set
%\[
%T(M) \defeq \fib\big(M_{h\cyc_2} \xrightarrow{\mathrm{Nm}} M^{h\cyc_2}\big).
%\]
%Thus \(T_2V=T(V^{\ot 2})\).
%The Tate spectral sequence for \(M\) is the homological spectral sequence
%\[
%E^2_{p,q}(M)\cong \widehat H_p(\cyc_2;H_q(M)) \Longrightarrow H_{p+q}(T(M)),
%\]
%where \(\widehat H_p\) denotes homological Tate group homology, normalized so that
%\[
%H_p(T(\Ftwo))\cong \widehat H_p(\cyc_2;\Ftwo).
%\]
%With this convention, over \(\Ftwo\),
%\[
%\widehat H_*(\cyc_2;\Ftwo)\cong \Ftwo[u,u^{-1}],
%\qquad |u|=1.
%\]
%We write
%\[
%\hat x_j \in \widehat H_j(\cyc_2;\Ftwo)
%\]
%for the generator satisfying
%\[
%u\,\hat x_j=\hat x_{j+1}.
%\]
%
%For \(M=V^{\ot 2}\), the Kunneth isomorphism gives
%\[
%H_q(V^{\ot 2})\cong \bigoplus_{a+b=q} H_a(V)\ot H_b(V),
%\]
%with \(\cyc_2\) acting by transposition.
%If \(e_\alpha,e_\beta\) are homogeneous classes in \(H(V)\), then the span
%\[
%\Ftwo\set{e_\alpha\ot e_\beta,e_\beta\ot e_\alpha}
%\]
%is induced from the trivial subgroup when \(\alpha\neq \beta\).
%Its Tate groups vanish.
%Thus only the diagonal classes
%\[
%e_\alpha\ot e_\alpha
%\]
%contribute to the \(E^2\)-page.
%
%We shall use the following convergence hypothesis.
%We say that \(V\) is \defn{Tate-admissible} if the Tate spectral sequence
%\[
%E^2_{p,q}(V^{\ot 2})\cong \widehat H_p(\cyc_2;H_q(V^{\ot 2})) \Longrightarrow H_{p+q}(T_2V)
%\]
%is strongly convergent, the diagonal classes
%\[
%\hat x_p\ot e_\alpha\ot e_\alpha
%\]
%survive to \(E^\infty\), and there are no hidden extensions among these classes as a graded \(\Ftwo[u,u^{-1}]\)-module.
%Equivalently, the associated graded description coming from the Tate spectral sequence lifts to an isomorphism of graded \(\Ftwo[u,u^{-1}]\)-modules
%\[
%\bigoplus_\alpha \Ftwo[u,u^{-1}]\set{\hat x_{-d_\alpha}\ot e_\alpha\ot e_\alpha}
%\xra{\cong} H(T_2V),
%\]
%where \(d_\alpha=|e_\alpha|\).
%
%This hypothesis is automatic, for example, if \(H(V)\) is finite-dimensional over \(\Ftwo\).
%Indeed, in that case \(V\) is isomorphic in \(\cD(\Ftwo)\) to a finite direct sum of shifts of \(\Ftwo\), and \(V^{\ot 2}\) decomposes as a finite direct sum of diagonal trivial \(\cyc_2\)-modules and off-diagonal induced \(\Ftwo[\cyc_2]\)-modules.
%The induced summands have zero Tate homology, and the diagonal summands contribute shifted copies of \(H(T(\Ftwo))\).
%
%\subsubsection{Tate diagonal}
%
%The \defn{Tate diagonal} is the natural transformation
%\[
%\Delta^t \colon \id_{\cD(\Ftwo)} \to T_2(-)
%\]
%which promotes the assignment \(v \mapsto v \ot v\) to a natural map in \(\cD(\Ftwo)\).
%On homology it is characterized by
%\[
%\begin{tikzcd}[row sep=0, column sep=small]
%	\Delta^t_* &[-20pt] \colon H(V) \arrow[r] & H(T_2V) \\
%	& v \arrow[r, |->] & \hat{x}_{-|v|} \ot v \ot v.
%\end{tikzcd}
%\]
%The map \(\Delta^t_*\) is linear because, for cycles \(v,w\), the cross term
%\[
%v\ot w+w\ot v
%\]
%is a norm class, hence maps to zero in Tate homology.
%
%\begin{lemma}\label{p:tate_iso}
%	Assume that \(V\) is Tate-admissible.
%	There is a natural isomorphism of graded \(\Ftwo[u,u^{-1}]\)-modules
%	\[
%	\Phi_V \colon \Ftwo[u,u^{-1}] \ot H(V) \xra{\cong\,} H(T_2V)
%	\]
%	given by
%	\[
%	\begin{tikzcd}[row sep=0, column sep=small]
%		\Ftwo[u,u^{-1}] \ot H(V) \arrow[r] & H(T_2V) \\
%		\hspace*{13pt} u^k \ot v \arrow[r, |->] & u^k \Delta^t_*(v).
%	\end{tikzcd}
%	\]
%\end{lemma}
%
%\begin{proof}
%	Choose a homogeneous basis \(\set{e_\alpha}_{\alpha\in A}\) of \(H(V)\), and write
%	\[
%	d_\alpha=|e_\alpha|.
%	\]
%	By the Kunneth isomorphism,
%	\[
%	H(V^{\ot 2})\cong H(V)^{\ot 2}.
%	\]
%	As a \(\cyc_2\)-module, \(H(V)^{\ot 2}\) decomposes into the diagonal summands
%	\[
%	\Ftwo\set{e_\alpha\ot e_\alpha}
%	\]
%	and the off-diagonal summands
%	\[
%	\Ftwo\set{e_\alpha\ot e_\beta,e_\beta\ot e_\alpha},
%	\qquad \alpha<\beta.
%	\]
%	Each off-diagonal summand is isomorphic to the induced module
%	\[
%	\operatorname{Ind}_1^{\cyc_2}\Ftwo.
%	\]
%	Induced modules have trivial Tate groups, so the off-diagonal summands contribute nothing to the \(E^2\)-page of the Tate spectral sequence.
%
%	The diagonal summand
%	\[
%	\Ftwo\set{e_\alpha\ot e_\alpha}
%	\]
%	is a copy of the trivial \(\cyc_2\)-module, placed in homological degree \(2d_\alpha\).
%	Its contribution to the \(E^2\)-page is therefore
%	\[
%	\widehat H_p(\cyc_2;\Ftwo)\set{e_\alpha\ot e_\alpha}
%	\]
%	in bidegree \((p,2d_\alpha)\).
%	The corresponding class in total degree \(p+2d_\alpha\) is
%	\[
%	\hat x_p\ot e_\alpha\ot e_\alpha.
%	\]
%	In particular, the unique Tate translate of this diagonal class lying in total degree \(d_\alpha\) is
%	\[
%	\hat x_{-d_\alpha}\ot e_\alpha\ot e_\alpha.
%	\]
%
%	By the defining property of the Tate diagonal,
%	\[
%	\Delta^t_*(e_\alpha)=\hat x_{-d_\alpha}\ot e_\alpha\ot e_\alpha.
%	\]
%	Multiplication by \(u\) gives
%	\[
%	u^k\Delta^t_*(e_\alpha)=\hat x_{k-d_\alpha}\ot e_\alpha\ot e_\alpha.
%	\]
%	These are exactly the diagonal Tate classes contributed by the basis vector \(e_\alpha\).
%
%	Since \(V\) is Tate-admissible, the Tate spectral sequence converges strongly, all these diagonal classes survive, and there are no hidden extensions as a graded \(\Ftwo[u,u^{-1}]\)-module.
%	Therefore the classes
%	\[
%	u^k\Delta^t_*(e_\alpha),
%	\qquad k\in \Z,\ \alpha\in A,
%	\]
%	form a homogeneous \(\Ftwo\)-basis of \(H(T_2V)\), compatible with the \(\Ftwo[u,u^{-1}]\)-module structure.
%	It follows that
%	\[
%	\Phi_V \colon \Ftwo[u,u^{-1}]\ot H(V) \to H(T_2V)
%	\]
%	is an isomorphism of graded \(\Ftwo[u,u^{-1}]\)-modules.
%
%	It remains only to check that the displayed formula is independent of the chosen homogeneous basis.
%	If \(v,w\in H(V)\) are homogeneous of the same degree, then
%	\[
%	(v+w)\ot(v+w)=v\ot v+w\ot w+v\ot w+w\ot v.
%	\]
%	The cross term \(v\ot w+w\ot v\) is a norm class.
%	Hence it maps to zero in Tate homology.
%	Therefore
%	\[
%	\Delta^t_*(v+w)=\Delta^t_*(v)+\Delta^t_*(w).
%	\]
%	The same argument, applied degreewise, gives additivity for arbitrary homogeneous decompositions.
%	Thus the formula
%	\[
%	v\mapsto \hat x_{-|v|}\ot v\ot v
%	\]
%	defines the same map independently of the chosen homogeneous basis.
%\end{proof}
%
%\subsubsection{Tate multiplications and Steenrod squares}
%
%A \defn{Tate multiplication} on \(V\) is a map
%\[
%\widehat\mu \colon T_2 V \to V.
%\]
%Its \(k^\th\) \defn{Steenrod square} is defined by composing the universal class \(u^k\Delta^t_*\) with \(\widehat\mu_*\).
%Explicitly,
%\[
%\Sq_k^{\widehat\mu} \colon H(V) \xra{\Delta^t_*} H(T_2 V) \xra{u^k} H(T_2 V) \xra{\widehat\mu_*} H(V).
%\]
%We say that \(\widehat\mu\) has:
%\begin{enumerate}
%	\item The \textit{identity property} if \(\Sq_0^{\widehat\mu} = \id_{H(V)}\),
%	\item The \textit{instability property} if \(\Sq_k^{\widehat\mu}(v) = 0\) for all \(v \in H(V)\) and \(k < |v|\),
%	\item The \textit{connectivity property} if \(\Sq_k^{\widehat\mu} = 0\) for all \(k < 0\).
%\end{enumerate}
%
%\subsubsection{Identity and instability}
%
%Consider a Tate multiplication \(\widehat\mu\) on \(V\).
%
%\begin{lemma}\label{l:identity_property}
%	\(\widehat\mu\) has the identity property if and only if
%	\[
%	\widehat\mu \circ \Delta^t \sim \id_V.
%	\]
%\end{lemma}
%
%\begin{proof}
%	By definition,
%	\[
%	\Sq_0^{\widehat\mu}=\widehat\mu_* \circ \Delta^t_*=(\widehat\mu\circ\Delta^t)_*.
%	\]
%	Thus \(\widehat\mu\) has the identity property if and only if
%	\[
%	(\widehat\mu\circ\Delta^t)_*=(\id_V)_*.
%	\]
%	Over a field, the derived category of chain complexes is equivalent to the category of graded vector spaces via homology.
%	Consequently, a morphism in \(\cD(\Ftwo)\) is determined by its induced map on homology.
%	Therefore the last equality holds if and only if
%	\[
%	\widehat\mu\circ\Delta^t \sim \id_V.
%	\]
%\end{proof}
%
%\begin{lemma}\label{l:instability_factorization}
%	\(\widehat\mu\) has the instability property if and only if it factors through homotopy orbits:
%	\[
%	\begin{tikzcd}[column sep=24pt]
%		(V^{\ot 2})^{h\cyc_2}[-1] \rar & T_2(V) \arrow[r] \arrow[d, "\widehat\mu"'] &
%		(V^{\ot 2})_{h\cyc_2} \arrow[dl, dashed, "\mu"] \\
%		& V. &
%	\end{tikzcd}
%	\]
%\end{lemma}
%
%\begin{proof}
%	Applying \([-,V]_{\cD(\Ftwo)}\) to the distinguished triangle \cref{eq:tate_sequence} gives an exact sequence.
%	Hence \(\widehat\mu\) factors through the projection
%	\[
%	T_2V \to (V^{\ot 2})_{h\cyc_2}
%	\]
%	if and only if the composite
%	\[
%	(V^{\ot 2})^{h\cyc_2}[-1] \to T_2V \xra{\widehat\mu} V
%	\]
%	is null.
%
%	Since morphisms in \(\cD(\Ftwo)\) are detected on homology, this composite is null if and only if \(\widehat\mu_*\) vanishes on the image of
%	\[
%	H\big((V^{\ot 2})^{h\cyc_2}[-1]\big) \to H(T_2V).
%	\]
%	Under the isomorphism of \cref{p:tate_iso}, this image is the subspace spanned by the classes
%	\[
%	u^k\Delta^t_*(v)=\hat x_{k-|v|}\ot v\ot v
%	\]
%	with
%	\[
%	k-|v|<0.
%	\]
%	Equivalently, it is spanned by the classes \(u^k\Delta^t_*(v)\) with \(k<|v|\).
%
%	Therefore the composite is null if and only if
%	\[
%	\widehat\mu_*\big(u^k\Delta^t_*(v)\big)=0
%	\]
%	for all homogeneous \(v\in H(V)\) and all \(k<|v|\).
%	By definition, this condition is precisely
%	\[
%	\Sq_k^{\widehat\mu}(v)=0
%	\]
%	for all \(v\) and all \(k<|v|\).
%\end{proof}
%
%\subsubsection{Coefficient sequence and connectivity}
%
%Consider the shifted Tate sequence for the coefficient module \(\Ftwo\), tensored with \(V\), where \(\Ftwo^{\ot 2}\cong \Ftwo\) carries the trivial \(\cyc_2\)-action:
%\[
%(\Ftwo^{\ot 2})^{h\cyc_2}[-1]\ot V \to T_2(\Ftwo)\ot V \to (\Ftwo^{\ot 2})_{h\cyc_2}\ot V.
%\]
%On homology this gives the coefficient sequence
%\[
%u^{-1}\Ftwo[u^{-1}]\ot H(V) \to \Ftwo[u,u^{-1}]\ot H(V) \to \Ftwo[u]\ot H(V).
%\]
%Here the first term is the negative Tate part, the middle term is the full Tate coefficient module, and the last term is the homotopy-orbit coefficient module.
%Equivalently, the map
%\[
%\Ftwo[u,u^{-1}]\ot H(V) \to \Ftwo[u]\ot H(V)
%\]
%is the quotient by the submodule
%\[
%u^{-1}\Ftwo[u^{-1}]\ot H(V).
%\]
%
%The isomorphism of \cref{p:tate_iso} identifies
%\[
%H(T_2V)
%\]
%with the middle term
%\[
%\Ftwo[u,u^{-1}]\ot H(V).
%\]
%This is a homological identification.
%We do not assert a canonical equivalence of functors
%\[
%T_2(V) \simeq T_2(\Ftwo)\ot V,
%\]
%since \(T_2(-)\) is a quadratic functor whereas \(T_2(\Ftwo)\ot -\) is linear.
%
%\begin{lemma}\label{l:connectivity_factorization}
%	\(\widehat\mu\) has the connectivity property if and only if the homology map
%	\[
%	\widehat\mu_* \colon H(T_2V) \to H(V)
%	\]
%	factors through the coefficient quotient
%	\[
%	\Ftwo[u,u^{-1}]\ot H(V) \to \Ftwo[u]\ot H(V)
%	\]
%	under the isomorphism of \cref{p:tate_iso}.
%	Equivalently, there exists a map
%	\[
%	\overline\mu_* \colon \Ftwo[u]\ot H(V) \to H(V)
%	\]
%	making the diagram
%	\[
%	\begin{tikzcd}[column sep=42pt]
%		\Ftwo[u,u^{-1}]\ot H(V) \arrow[r] \arrow[dr, "\widehat\mu_* \circ \Phi_V"'] &
%		\Ftwo[u]\ot H(V) \arrow[d, dashed, "\overline\mu_*"] \\
%		& H(V)
%	\end{tikzcd}
%	\]
%	commute.
%\end{lemma}
%
%\begin{proof}
%	The kernel of the quotient
%	\[
%	\Ftwo[u,u^{-1}]\ot H(V) \to \Ftwo[u]\ot H(V)
%	\]
%	is
%	\[
%	u^{-1}\Ftwo[u^{-1}]\ot H(V).
%	\]
%	Under the isomorphism
%	\[
%	\Phi_V \colon \Ftwo[u,u^{-1}]\ot H(V) \xra{\cong} H(T_2V),
%	\]
%	this kernel corresponds to the subspace spanned by the classes
%	\[
%	u^k\Delta^t_*(v)
%	\]
%	with \(k<0\).
%
%	Therefore \(\widehat\mu_*\) factors through the quotient if and only if
%	\[
%	\widehat\mu_*\big(u^k\Delta^t_*(v)\big)=0
%	\]
%	for every homogeneous \(v\in H(V)\) and every \(k<0\).
%	By definition, this condition is exactly
%	\[
%	\Sq_k^{\widehat\mu}(v)=0
%	\]
%	for every \(v\in H(V)\) and every \(k<0\).
%	Thus the factorization condition is equivalent to the connectivity property.
%\end{proof}
%
%
%\newpage


\section{Steenrod squares}

We work in the derived category \(\cD(\Ftwo)\) of unbounded chain complexes of vector spaces over the field with two elements.
Throughout this subsection we might need to assume that \(V\) is of finite type, that is, degreewise finite-dimensional.\anibal{Check this for the first lemma.}

\subsubsection{Shifted Tate sequence}

We use the symmetry of the tensor product to define the functors
\[
V \mapsto (V^{\ot 2})_{h\cyc_2}
\qquad\text{and}\qquad
V \mapsto (V^{\ot 2})^{h\cyc_2},
\]
where \((-)_{h\cyc_2}\) and \((-)^{h\cyc_2}\) are respectively the functors of homotopy orbits and homotopy fixed points.
The homotopy fibre of the norm map defines another functor
\[
T_2 V \defeq \fib\Big((V^{\ot 2})_{h\cyc_2} \xrightarrow{\mathrm{Nm}} (V^{\ot 2})^{h\cyc_2}\Big),
\]
naturally equivalent to the shifted Tate construction \(T_2 V \simeq (V^{\ot 2})^{t\cyc_2}[-1]\), and which fits in the \textit{shifted Tate sequence}
\begin{equation}\label{eq:tate_sequence}
	(V^{\ot 2})^{h\cyc_2}[-1] \to T_2 V \to (V^{\ot 2})_{h\cyc_2}.
\end{equation}

As is well-known, for every \(j \in \Z\),
\[
H_j(T_2 \Ftwo) \cong \Ftwo\set{\hat x_j},
\]
so \(H(T_2 \Ftwo)\) is a free one-dimensional module over the ring \(\Ftwo[u,u^{-1}]\) acting by
\[
u \, \hat x_j = \hat x_{j + 1}.
\]

\subsubsection{Tate diagonal}
\begin{definition}
The \defn{Tate diagonal} is the natural transformation
\[
\Delta^t \colon \id_{\cD(\Ftwo)} \to T_2(-)
\]
whose existence promotes the non-linear assignment \(v \mapsto v \ot v\) to an actual map in \(\cD(\Ftwo)\), and which is characterized on homology by
\[
\begin{tikzcd}[row sep=0, column sep=small]
	\Delta^t_* &[-20pt] \colon H(V) \arrow[r] & H(T_2V) \\
	& v \arrow[r, |->] & \hat{x}_{-|v|} \ot v \ot v.
\end{tikzcd}
\]
\end{definition}
The map \(\Delta^t_*\) is linear since the cross term \(v \ot w + w \ot v\) vanishes as it is a norm class. \martin{This is true for $p=2$, but I think for odd $p$ one needs to take the cofibre, not the fibre, because the argument will use that $(1-p)|v|$ is even, and so $d \hat x_{(1-p)|v|+1}=N \hat x_{(1-p)|v|}$, which only happens for odd degrees (even degrees have $T$).}

\medskip
The image of \(\Delta^t_*\) generates \(H(T_2 V)\) as a free \(\Ftwo[u,u^{-1}]\)-module with action
\[
u \Delta^t_*(v) = u (\hat x_{-|v|} \ot v \ot v) = \hat x_{1-|v|} \ot v \ot v.
\]
In other words, we have the following.

\begin{lemma}\label{p:tate_iso}
	The following map is an isomorphism of free \(\Ftwo[u,u^{-1}]\)-modules:
	\[
	\begin{tikzcd}[row sep=0, column sep=small]
		\Ftwo[u,u^{-1}] \ot H(V) \arrow[r] & H(T_2V) \\
		\hspace*{13pt} u^k \ot v \arrow[r, |->] & u^k \Delta^t_*(v).
	\end{tikzcd}
	\]
\end{lemma}

\subsubsection{Tate multiplications and Steenrod squares}

A \defn{Tate multiplication} on \(V\) is a map
\[
\widehat\mu \colon T_2 V \to V.
\]
Its \(k^\th\) \defn{Steenrod square} is defined by composing the \defn{universal \(k^\th\) Steenrod square} \(u^k \Delta^t_*\) with \(\widehat\mu_*\).
Explicitly,
\[
\Sq_k^{\widehat\mu} \colon H(V) \xra{\Delta^t_*} H(T_2 V) \xra{u^k} H(T_2 V) \xra{\widehat\mu_*} H(V).
\]
We say that \(\widehat\mu\) has:
\begin{enumerate}
	\item The \textit{identity property} if \(\Sq_0^{\widehat\mu} = \id_{H(V)}\),
	\item The \textit{instability property} if \(\Sq_k^{\widehat\mu}(v) = 0\) for all \(v \in H(V)\) and \(k < |v|\),
	\item The \textit{connectivity property} if \(\Sq_k^{\widehat\mu} = 0\) for all \(k < 0\).
\end{enumerate}
Observe that the instability property implies the connectivity one, since any $k<0$ will satisfy $k<|v|$ for any $v\in H(V)$.
\subsubsection{Identity and instability}

Consider a Tate multiplication \(\widehat\mu\) on \(V\).

\begin{lemma}
	\(\widehat\mu\) has the identity property if and only if
	\[
	\widehat\mu \circ \Delta^t \sim \id_V.
	\]
\end{lemma}

\begin{lemma}
	\(\widehat\mu\) has the instability property if and only if it factors as follows:
	\[
	\begin{tikzcd}[column sep=24pt]
		(V^{\ot 2})^{h\cyc_2}[-1] \rar & T_2(V) \arrow[r] \arrow[d, "\widehat\mu"'] &
		(V^{\ot 2})_{h\cyc_2} \arrow[dl, dashed, "\mu"] \\
		& V. &
	\end{tikzcd}
	\]
\end{lemma}

\subsubsection{Coefficients sequence and connectivity}

Consider the shifted Tate sequence for the coefficient module \(\Ftwo\), tensored with \(V\), where \(\Ftwo^{\ot 2} \cong \Ftwo\) carries the trivial \(\cyc_2\)-action:
\[
(\Ftwo^{\ot 2})^{h\cyc_2}[-1] \ot V \to T_2(\Ftwo) \ot V \to (\Ftwo^{\ot 2})_{h\cyc_2} \ot V.
\]
On homology this sequence identifies with
\[
u^{-1}\Ftwo[u^{-1}] \ot H(V) \to \Ftwo[u,u^{-1}] \ot H(V) \to \Ftwo[u] \ot H(V)
\]
or, more sugestively,
\[
H(V; H^{\bullet-1}(\cyc_2)) \to H(V; \widehat H_{\bullet}(\cyc_2)) \to H(V; H_{\bullet}(\cyc_2)).
\]
Recall from \cref{p:tate_iso} that the middle term is isomorphic to \(H(T_2 V)\).

\begin{lemma}
	A Tate multiplication \(\widehat\mu\) on \(V\) has the connectivity property if and only if it factors as follows:
	\[
	\begin{tikzcd}[column sep=24pt]
		(\Ftwo^{\ot 2})^{h\cyc_2}[-1] \ot V \rar &
		T_2(\Ftwo) \ot V \arrow[r] &
		(\Ftwo^{\ot 2})_{h\cyc_2} \ot V. \arrow[dashed, ddl]\\
		& T_2(V) \dar["\widehat\mu"] \uar["\sim"]& \\
		& V. &
	\end{tikzcd}
	\]
\end{lemma}